\documentclass[10pt,a4paper]{amsart}
\usepackage[margin=19mm]{geometry}
\usepackage{amsmath,amssymb,mathtools}
\usepackage[scr=rsfs,cal=euler]{mathalfa}
\usepackage{xfrac}
\usepackage[unicode,hidelinks,bookmarks=false]{hyperref}
\newcommand{\ie}{i.e.\ }
\newcommand{\cf}{cf.\ }
\newcommand{\resp}{respectively\ }
\newcommand{\Subsection}{\subsection}
\providecommand{\nobreakdash}{\nobreak\hbox{-}}
\setlength{\emergencystretch}{2em}
\multlinegap0pt
\usepackage{bm}
\usepackage{bbm}
\usepackage{dsfont}
\usepackage{xspace}
\usepackage{cancel}
\usepackage{mathtools}
\usepackage{amsthm}
\makeatletter
\@ifundefined{corollary}{\newtheorem{corollary}{Corollary}}{}
\@ifundefined{lemma}{\newtheorem{lemma}{Lemma}}{}
\makeatother
\def\R{{\mathbb R}}
\def\Z{{\mathbb Z}}
\def\cL{{\mathcal L}}
\def\cM{{\mathcal M}}
\def\cR{{\mathcal R}}
\def\ge{\geqslant}
\def\le{\leqslant}
\title[Weyl sums with multiplicative coefficients and joint equidis]{Weyl sums with multiplicative coefficients and joint equidistribution}
\author{Matteo Bordignon, Cynthia Bortolotto,, Bryce Kerr}
\date{Source: arXiv:2303.03768v1 (CC BY 4.0)}
\begin{document}
\maketitle

\noindent\emph{Source and licence.} Weyl sums with multiplicative coefficients and joint equidistribution. Version arXiv:2303.03768v1 (CC BY 4.0), distributed under the Creative Commons Attribution 4.0 International licence, as recorded in the arXiv licence metadata for this exact version and shown on the abstract page https://arxiv.org/abs/2303.03768v1. Full rights notice: licenses/arxiv-2303.03768-CC-BY-4.0.txt.\par\smallskip

\noindent\textit{Editorial scope.} Counted unit: the Lemma whose source label is \texttt{lem:bilinear1} in the article (source0.tex lines 643--697, source label \texttt{lem:bilinear1}), delivered together with the article's own complete original proof of it (source0.tex lines 699--856, one \texttt{proof} environment of 158 lines). No mathematics is written, repaired or re-derived: statement and proof are the article's own text line for line. Required context retained uncounted: eq:cor76867 (lines 595--608); lem:vmvtP1 (lines 623--633). Every definition, display and label of those lines is retained unchanged. Omitted from this package and not redistributed: 1--594, 609--622, 634--642, 698--698, 857--1710. Nothing inside a retained excerpt is omitted or altered: every retained body is delivered exactly as the article's lines are written, so no omission receipt of any kind accompanies this package. Citations: the retained excerpts carry no citation command at all, so no bibliography entry of the article is transcribed and no reference is redistributed. Reference adaptation: none was needed and none was made; every reference inside the delivered unit resolves inside it, so no unresolved reference or citation is produced. Printed slips kept literal: no printed word, symbol, hypothesis or number of the counted statement or of its proof was altered. Layout and class: the article's own class is \texttt{amsart} with packages including amsmath, amssymb, amsbsy, amsfonts, amsthm, latexsym, amsopn, amstext, amsxtra, euscript, amscd, mathrsfs, color, bm; the wrapper is the lane's pinned \texttt{amsart} editorial wrapper at 10pt on A4, which restates the theorem-like environments the retained text needs and adds the article's own macro definitions that the retained text needs (\texttt{\textbackslash R}, \texttt{\textbackslash Z}, \texttt{\textbackslash cL}, \texttt{\textbackslash cM}, \texttt{\textbackslash cR}, \texttt{\textbackslash ge}, \texttt{\textbackslash le}), with extra wrapper packages (bm, bbm, dsfont, xspace, cancel, mathtools). The delivered numbering is the unit's own. Assets: no retained span contains a figure environment, a graphics-inclusion command, a PDF-page inclusion, a table or a TikZ picture; no image file is copied into this package, the package declares no asset dependency and no image file is redistributed with it. Licence: version arXiv:2303.03768v1 of this article is distributed under the Creative Commons Attribution 4.0 International licence, as recorded in the arXiv licence metadata for this exact version and shown on the abstract page https://arxiv.org/abs/2303.03768v1. A full rights notice is delivered at licenses/arxiv-2303.03768-CC-BY-4.0.txt. Characters: every retained excerpt is pure ASCII, so the delivered engine, XeLaTeX with its default Latin Modern text font, reports no missing character.
\begin{corollary}
\label{eq:cor76867}
Let $d\ge 2$, $r>d(d+1)$ be integers and $\cM(k), 1\le k \le K$ disjoint intervals satisfying 
\begin{align*}
\cM(k)=(M'(k),M''(k)], \quad M''(k)-M'(k)\le Y.
\end{align*}
Then the number of solutions to the system of equations 
\begin{align}
\label{eq:system1}
n_1^{j}+\dots-n_{2r}^{j}=0, \quad n_i\in \cM(k), \quad 1\le k \le K,
\end{align}
is bounded by 
$$O(KY^{2r-d(d+1)/2}).$$
\end{corollary}

\begin{lemma}
\label{lem:vmvtP1}
Let $d\ge 2$ be an integer and $X, Y\gg 1$. If $r>d(d+1)$,  the number of solutions to the equation 
\begin{align*}
&p_1^{j}+\dots-p_{2r}^{j}=0, \quad 1\le j \le d, \quad Y\le p_i\le Y+X, \\ & X\gg\frac{Y}{(\log Y)^M}, \quad p_i \quad \text{prime},
\end{align*}
for any $M\gg1$, is bounded by 
\begin{align*}
O\left(\frac{X^{2r-d(d+1)/2}}{(\log{X})^{2r}}\right).
\end{align*}
\end{lemma}

\begin{lemma}
\label{lem:bilinear1}
Let $F\in \R[X]$ be a polynomial of degree $d\ge 2$  of the form 
$$F(x)=\alpha_d x^{d}+\dots +\alpha_1 x.$$
 Let $M,N$ be integers and $\alpha(n),\beta(m)$ two sequences of complex numbers satisfying 
\begin{align*}
|\beta(m)|\le 1,
\end{align*}
and
\begin{align}
\label{eq:condi}
\sum_{n\le N} |\alpha (n)| \ll N, \quad \sum_{n\le N} |\alpha (n)|^2 \ll N(\log{N})^A,
\end{align}
for $A \ge 0$.
For $1\le k \le K$ let 
\begin{align}
\label{eq:Rk}
\cR(k)=\cL(k)\times \cM(k),
\end{align}
 be a rectangle of the form 
$$\cL(k)=(Q'(k),Q''(k)], \quad \cM(k)=(M'(k)\times M''(k)],$$
such that $\cL(k)\subseteq (0,Q)$, with $Q\gg 1$, are disjoint and satisfy 
\begin{equation}
\label{eq:cond1}
Q''(k)-Q'(k)\le X,
\end{equation}
or
\begin{equation}
\label{eq:cond2}
Q''(k)-Q'(k)= X \gg \frac{ Q'(k)}{(\log Q'(k))^M}
\end{equation}
and $\cM(k)\subseteq (0,M]$ are disjoint and satisfy 
$$M''(k)-M'(k)\le Y, \quad M''(k)\le 2M'(k),$$
and $K\ll M$.
Let $I$ denote the sum 
\begin{align}
\label{eq:Idef}
I=\sum_{k=1}^{K}\sum_{\substack{(p,n)\in \cR(k)}}\alpha(n)\beta(p)e(F(np)).
\end{align}
Let $R\ge 1$ and suppose there exists $q\le R$ and $1\le \ell \le d$ such that 
\begin{align}
\label{eq:alassump}
\left|\alpha_{\ell}-\frac{a}{q}\right|\le \frac{1}{qR},
\end{align}
for some $(a,q)=1$. For $r>d(d+1)$, we have

\begin{align*}
I^{4r^2}& \ll m(X)(\log{M})^{2rA}M^{4r^2}\left(\frac{X}{\log{X}}\right)^{4r^2} \left(\frac{M}{Y}\right)^{d(d-1)/2}\\ &\left(\frac{Q}{X}\right)^{d(d-1)/2}  \left(\frac{q}{XQ^{\ell-1}YM^{\ell-1}}+\frac{1}{XQ^{\ell-1}}+\frac{1}{YM^{\ell-1}}+\frac{1}{q}\right),
\end{align*}
where
$$m(x)=\begin{cases*}
(\log{x})^{4r}& \textit{when} \eqref{eq:cond1} \textit{holds},\\
1 & \textit{when} \eqref{eq:cond2} \textit{holds}.
\end{cases*} $$
\end{lemma}

\begin{proof}

Recall that
$$F(x)=\alpha_d x^d+\dots+\alpha_1 x,$$
and define
\begin{align}
\label{eq:Sums1}
S_k=\sum_{\substack{(p,n)\in \cR(k)}}\alpha(n)\beta(p)e(F(np)),
\end{align}
so that 
\begin{align}
\label{eq:IS}
I=\sum_{k=1}^{K}S_k.
\end{align}
Fix some $1\le k \le K$ and consider~\eqref{eq:Sums1}. Recalling~\eqref{eq:Rk}, by H\"{o}lder's inequality, for any integer $r\ge 1$
\begin{align*}
&S_{k}\le \\ & \left(\sum_{n\in \cM(k)}|\alpha(n)|^{2r/(2r-1)}\right)^{1-1/2r}\left(\sum_{n\in \cM(k)}\left|\sum_{p\in \cL(k)}\beta(p)e(F(np))\right|^{2r}\right)^{1/2r}.
\end{align*}
After another application of H\"{o}lder's inequality 
\begin{align*}
&I^{2r}\le \\
&\left(\sum_{1\le k \le K}\sum_{n\in \cM(k)}|\alpha(n)|^{2r/(2r-1)} \right)^{2r-1}\sum_{1\le k \le K}\sum_{n\in \cM(k)}\left|\sum_{p\in \cL(k)}\beta(p)e(F(np))\right|^{2r}.
\end{align*}
Hence from assumptions on the $\cM(k)$
\begin{align*}
I^{2r}\ll \left( \sum_{n \le M} |\alpha (n)|^{2r/(2r-1)} \right)^{2r-1}\sum_{1\le k \le K}\sum_{n\in \cM(k)}\left|\sum_{p\in \cL(k)}\beta(p)e(F(np))\right|^{2r},
\end{align*}
which combined with~\eqref{eq:condi} implies that 
\begin{align*}
I^{2r}\ll  M^{2r-1}(\log{M})^{A } \sum_{1\le k \le K}\sum_{n\in \cM(k)}\left|\sum_{p\in \cL(k)}\beta(p)e(F(np))\right|^{2r}.
\end{align*}
Write $\lambda=(\lambda_1,\dots,\lambda_d)$ and let $J_k(\lambda)$ count the number of solutions to 
$$p_1^{j}+\dots-p_{2r}^j=\lambda_j, \quad 1\le j \le d, \quad p_i\in \cL(k), \quad p_i \ \ \text{prime}.$$
Note by assumptions on $\cL(k)$, if $p_1,\dots,p_{2r}\in \cL(k)$ satisfy 
$$p_1^{j}+\dots-p_{2r}^j=\lambda_j,$$
then 
$$\lambda_j\ll XQ^{j-1}.$$
Expanding the $2r$-th power and interchanging summation gives
\begin{align*}
I^{2r}\ll& M^{2r-1}(\log{M})^{A}\\ & \times \sum_{1\le k \le K}\sum_{\substack{\lambda \\ \lambda_j \ll XQ^{j-1}}} J_k(\lambda)\left|\sum_{n\in \cM(k)}e(\alpha_d\lambda_d n^{d}+\dots+\alpha_1\lambda_1 n)\right|.
\end{align*}
After another application of H\"{o}lder's inequality, we get  
\begin{align*}
& I^{4r^2}\ll M^{4r^2-2r}(\log{M})^{2rA}\left(\sum_{1\le k \le K}\sum_{\lambda} J_k(\lambda)\right)^{2r-2}\left(\sum_{1\le k \le K}\sum_{\lambda} J_k(\lambda)^2\right) \\ 
& \quad \times \sum_{1\le k \le K}\sum_{\substack{ \lambda_j \ll XQ^{j-1}}}\left|\sum_{n\in \cM(k)}e(\alpha_d\lambda_d n^{d}+\dots+\alpha_1\lambda_1 n)\right|^{2r}.
\end{align*}

We have 
\begin{align*}
\sum_{1\le k \le K}\sum_{\lambda} J_k(\lambda)\le \sum_{1\le k \le K}|\{ p_1,\dots,p_{2r}\in \cL(k)\}|.
\end{align*}



Since for each $1\le k \le K$,  $\cL(k)$ is an interval of length at most $X$, the Brun-Titshmarsh theorem implies 
\begin{align*}
\sum_{1\le k \le K}\sum_{\lambda} J_k(\lambda)\ll \frac{KX^{2r}}{(\log{X})^{2r}}.
\end{align*}
The term
\begin{align*}
\sum_{1\le k \le K}\sum_{\lambda} J_k(\lambda)^2,
\end{align*}
counts the number of solutions to the system of equations 
\begin{align}
\label{eq:123456-1}
p_1^{j}+\dots-p_{4r}^{j}=0, \quad 1\le j \le d, \quad p_i\in \cL(k), \ \ \text{ $p_i$ prime}, \quad 1\le k \le K.
\end{align}
Ignoring the condition that $p_i$ is prime, by Corollary~\ref{eq:cor76867} and Lemma \ref{lem:vmvtP1},
\begin{align*}
\sum_{1\le k \le K}\sum_{\lambda} J_k(\lambda)^2\ll K m_1(X), 
\end{align*}
where
$$m_1(x)=\begin{cases*}
X^{4r-d(d+1)/2}& \textit{when} \eqref{eq:cond1} \textit{holds},\\
\frac{X^{4r-d(d+1)/2}}{(\log X)^{4r}} & \textit{when} \eqref{eq:cond2} \textit{holds}.
\end{cases*} $$
since $r>d(d+1)$. Combined with the observations above, this implies 
\begin{align}
\label{eq:II0}
& I^{4r^2}\ll m(X) M^{4r^2-2r}(\log{M})^{2rA}K^{2r-1}\left(\frac{X}{\log{X}}\right)^{4r^2}\frac{I_0}{X^{d(d+1)/2}},
\end{align}
where 
$$I_0=\sum_{1\le k \le K}\sum_{\substack{ \lambda_j \ll XQ^{j-1}}}\left|\sum_{n\in \cM(k)}e(\alpha_d\lambda_d n^{d}+\dots+\alpha_1\lambda_1 n)\right|^{2r}.$$
Let 
$$S(x)=\left(\frac{\sin{\pi x}}{\pi x}\right)^2,$$
so that 
\begin{align}
\label{eq:Shat}
\widehat S(x)=\max\{0,1-|x|\},
\end{align}
and 
$$S(x) \gg 1 \quad \text{if} \quad |x|\le \frac{1}{4}.$$
There exists a constant $c_0$ such that 
\begin{align*}
&I_0\ll \\ & \sum_{1\le k \le K}\sum_{\substack{ \lambda_j \ll XQ^{j-1} \\  j\neq \ell}}\sum_{\lambda_{\ell}\in \Z}S\left(\frac{c_0\lambda_{\ell}}{XQ^{\ell-1}}\right)\left|\sum_{n\in \cM(k)}e(\alpha_d\lambda_d n^{d}+\dots+\alpha_1\lambda_1 n)\right|^{2r}.
\end{align*}

Expanding the $2r$-th power and interchanging summation 
\begin{align*}
&I_0\ll \\ &\sum_{\substack{1\le k \le K\\ \mu_j\ll YN^{j-1}}}L_k(\mu)\prod_{\substack{j=1 \\ j\neq \ell }}^{d}\left|\sum_{\lambda\ll XQ^{j-1}}e(\alpha_j \lambda \mu_j)\right|\left|\sum_{\lambda \in \Z}S\left(\frac{c_0\lambda_{\ell}}{XQ^{\ell-1}}\right)e(\alpha_\ell \lambda \mu_\ell)\right|,
\end{align*}
with $\mu=(\mu_1,\dots,\mu_d)$ and $L_k(\mu)$ counts the number of solutions to 
\begin{align*}
n^{j}_1+\dots-n_{2r}^{j}=\mu_j, \quad n_i\in \cM(k), \quad 1\le k \le K.
\end{align*}
Using that 
$$L_k(\mu)\le L_k(0),$$
and applying Corollary~\ref{eq:cor76867}, we obtain

\begin{align*}
& I_0\ll KY^{2r-d(d+1)/2} \\ & \sum_{\substack{\mu \\ \mu_j \ll YM^{j-1}}}\prod_{\substack{j=1 \\ j\neq \ell }}^{d}\left|\sum_{\lambda \ll XQ^{j-1}}e(\alpha_j \lambda \mu_j)\right|\left|\sum_{\lambda \in \Z}S\left(\frac{c_0\lambda}{Xq^{\ell-1}}\right)e(\alpha_\ell \lambda \mu_\ell)\right|.
\end{align*}
Combining the inequality above with~\eqref{eq:II0} we have 
\begin{align}
\label{eq:II1}
& I^{4r^2}\ll m(X)(\log{M})^{2rA}M^{4r^2}\left(\frac{X}{\log{X}}\right)^{4r^2}\frac{I_1}{(XY)^{d(d+1)/2}},
\end{align}
where 

\begin{align*}
I_1=\sum_{\substack{\mu_j \ll YM^{j-1}}}\prod_{\substack{j=1 \\ j\neq \ell }}^{d}\left|\sum_{\lambda \ll XQ^{j-1}}e(\alpha_j \lambda \mu_j)\right|\left|\sum_{\lambda \in \Z}S\left(\frac{c_0\lambda}{Xq^{\ell-1}}\right)e(\alpha_\ell \lambda \mu_\ell)\right|
\end{align*}
In $I_1$, we bound every term trivially except the one with index $\ell$ to obtain 
\begin{align}
\label{eq:I1I1}
I_1&\ll \left(\frac{Y}{M}\right)^{d}\left(\frac{X}{Q}\right)^{d}(QM)^{d(d+1)/2}\frac{1}{XQ^{\ell-1}}\frac{1}{YM^{\ell-1}} \\ &\sum_{\mu\ll YM^{\ell-1}}\left|\sum_{\lambda \in \Z}S\left(\frac{c_0\lambda}{Xq^{\ell-1}}\right)e(\alpha_\ell \lambda \mu_\ell)\right|.
\end{align}
By Poisson summation
\begin{align*}
\sum_{\lambda \in \Z}S\left(\frac{c_0\lambda_{\ell}}{XQ^{\ell-1}}\right)e(\alpha_\ell\mu \lambda)=\frac{XQ^{\ell-1}}{c_0}\sum_{\lambda \in \Z}\widehat S\left(\frac{XQ^{\ell-1}(\lambda-\alpha_{\ell}\mu)}{c_0}\right),
\end{align*}
and hence from~\eqref{eq:Shat}
\begin{align*}
&\sum_{\mu\ll YM^{\ell-1}}\left|\sum_{\lambda \in \Z}S\left(\frac{c_0\lambda_{\ell}}{XQ^{\ell-1}}\right)e(\alpha_\ell\mu \lambda)\right| \\ &  \quad \quad \quad \quad \quad \ll XQ^{\ell-1}\sum_{\mu\ll YM^{\ell-1}}\sum_{\lambda \in \Z}\max\left\{0,1-\left|\frac{XQ^{\ell-1}(\lambda-\alpha_{\ell}\mu)}{c_0} \right| \right\} \\ & 
\quad \quad \quad \quad \quad \ll XQ^{\ell-1}\left|\left\{ \mu \ll YM^{\ell-1} \ : \ \|\alpha_{\ell}\mu\|\le \frac{100c_0}{XQ^{\ell-1}} \right\}\right|.
\end{align*}
There exists some real number $\beta$ such that 
\begin{align*}
&\left|\left\{ \mu \ll YM^{\ell-1} \ : \ \|\alpha_{\ell}\mu\|\le \frac{100c_0}{XQ^{\ell-1}} \right\}\right| \\ & \quad \quad \quad \ll \left(1+\frac{YM^{\ell-1}}{q}\right)\left|\left\{ 0\le \mu \le q \ : \ \|\alpha_{\ell}\mu+\beta\|\le \frac{100c_0}{XQ^{\ell-1}} \right\}\right|.
\end{align*}
Recalling~\eqref{eq:alassump}
\begin{align*}
&\left|\left\{ \mu \ll YM^{\ell-1} \ : \ \|\alpha_{\ell}\mu\|\le \frac{100c_0}{XQ^{\ell-1}} \right\}\right| \\ &  \quad \quad \ll \left(1+\frac{YM^{\ell-1}}{q}\right)\left|\left\{ 0\le \mu \le q \ : \ \|\frac{a\mu}{q}+\beta\|\le \frac{100c_0}{XQ^{\ell-1}}+\frac{1}{R} \right\}\right|,
\end{align*}
which implies that 
\begin{align*}
&\left|\left\{ \mu \ll YM^{\ell-1} \ : \ \|\alpha_{\ell}\mu\|  \le \frac{100c_0}{XQ^{\ell-1}} \right\}\right|\\ & \quad \quad \quad \quad \quad \quad \ll \left(1+\frac{YM^{\ell-1}}{q} \right)\left(1+\frac{q}{XQ^{\ell-1}}\right),
\end{align*}
and hence 
\begin{align*}
&\sum_{\mu\ll YM^{\ell-1}}\left|\sum_{\lambda \in \Z}S\left(\frac{c_0\lambda}{Xq^{\ell-1}}\right)e(\alpha_\ell \lambda \mu_\ell)\right|\ll \\ & XQ^{\ell-1}YM^{\ell-1}\left(\frac{1}{YM^{\ell-1}}+\frac{1}{q} \right)\left(1+\frac{q}{XQ^{\ell-1}}\right).
\end{align*}
Combined with~\eqref{eq:II1} and~\eqref{eq:I1I1} gives 
\begin{align*}
& I^{4r^2}\ll \\ & m(X)(\log{M})^{2rA}M^{4r^2}\left(\frac{X}{\log{X}}\right)^{4r^2} \left(\frac{M}{Y}\right)^{d(d-1)/2}\left(\frac{Q}{X}\right)^{d(d-1)/2} \\ & \left(\frac{q}{XQ^{\ell-1}YM^{\ell-1}}+\frac{1}{XQ^{\ell-1}}+\frac{1}{YM^{\ell-1}}+\frac{1}{q}\right),
\end{align*}
which completes the proof.
\end{proof}

\end{document}
